% !PRINTED 350
\weberpage{350}
%<<CONT
conjugirte Paare sein. Dann aber bilden die drei Wurzeln $(1,1)$ $(2,0)$ $(0,2)$ ein Tripel, deren conjugirte $(1,1)$ $(1,0)$ $(0,1)$ kein Tripel bilden, was unmöglich ist. Jede der drei reellen Wurzeln $(0,0)$, $(1,1)$, $(2,2)$ kommt also ausser in (1) noch in einem und nur in einem reellen Tripel vor, dessen beide anderen Wurzeln conjugirt imaginär sind, und der dritte Fall führt also zu der Anordnung der reellen und imaginären Wurzeln, die wir bei den Wendepunkten der Curve dritter Ordnung kennen gelernt haben.

Dass aber auch die Fälle 1. und 2. vorkommen können, lehrt folgende Betrachtung: Aus einer allgemeinen Gleichung $9^{\text{ten}}$ Grades erhält man eine Tripelgleichung durch Adjunction einer zu der linearen Gruppe $SQ$ gehörigen Function. Eine solche Function ist z.\,B.:

\[
\begin{aligned}
v &= (0,0)(1,1)(2,2) + (0,0)(1,2)(2,1) + (0,0)(1,0)(2,0) \\
&\quad + (0,0)(0,1)(0,2) + (1,1)(1,2)(1,0) + (1,1)(2,1)(0,1) \\
&\quad + (1,1)(0,2)(2,0) + (1,2)(0,1)(2,0) + (2,2)(2,1)(2,0) \\
&\quad + (2,2)(1,2)(0,2) + (2,2)(0,1)(1,0) + (2,1)(0,2)(1,0),
\end{aligned}
\]

und diese Function ist reell in den Fällen 1., 2., 3. Also ist auch der erweiterte Rationalitätsbereich, in dem unsere Gleichung eine Tripelgleichung ist, reell.

\begin{center}
\rule{0.22\textwidth}{0.4pt}
\end{center}
